\documentclass[11pt,a4paper]{article}

\usepackage[T1]{fontenc}
\usepackage[utf8]{inputenc}
\usepackage[margin=1.1in]{geometry}
\usepackage{amsmath,amssymb,amsthm}
\usepackage{booktabs}
\usepackage{natbib}
\usepackage{hyperref}
\hypersetup{colorlinks=true, linkcolor=black, citecolor=black, urlcolor=blue}
\usepackage{orcidlink}

\newtheorem{proposition}{Proposition}
\newtheorem{remark}{Remark}
\theoremstyle{definition}
\newtheorem{definition}{Definition}

\title{Fail-closed prediction: refusal as the boundary of a\\
distribution-free guarantee}

\author{%
  Diogo Ribeiro\,\orcidlink{0009-0001-2022-7072}\\[2pt]
  \small Faculty of Media Arts and Design, Technical University of Porto,
  Portugal
}

\date{Draft --- 22 September 2026}

\begin{document}
\maketitle

\begin{abstract}
A predictive model in general-purpose statistical software answers every
question put to it, including questions it has no evidence about. We argue
that the scalar prediction is the wrong return type, and that a prediction
should carry a value, a guarantee and a validity status, with refusal as an
admissible outcome. The argument is not a matter of engineering taste. Under
covariate shift, the weighted conformal construction of
\citet{tibshirani2019} assigns a point mass at infinity to the test point
itself; when that mass exceeds the miscoverage level the prediction set is the
whole outcome space. Refusal is the finite-sample expression of that event, and
returning \texttt{NA} with a reason is a more faithful report than returning an
unbounded interval or, worse, the narrow one that an unweighted quantile would
have produced. We then observe that the thresholds used to decide whether a
model is fit to deploy are routinely compared against statistics whose null
behaviour depends on sample size, class balance or estimation in sample, and
give four measured instances in which a conventional threshold misfires
severely: the population stability index flags roughly 80\% of unshifted
50-row batches; an in-sample classifier two-sample area under the curve reads
0.53 on identical samples; single-feature leakage screening by accuracy scores
0.974 on pure noise when the event rate is 2.6\%; and point-estimate
comparison fails checks on sampling noise alone. Each is repaired by
calibrating the threshold against a null rather than a convention. The
\textsf{R} package \texttt{attest} implements the resulting contract and is
the proof of concept.
\end{abstract}

\section{Introduction}

Fitting a model and validating it are separate acts, performed by the same
person minutes apart, and thereafter linked by nothing but memory. The fitted
object records its coefficients, not the leakage screen that preceded it; the
notebook records the calibration plot, not which object it described. This is
not a failure of diligence. It is a consequence of the return type: a model is
a function to a scalar, and a scalar cannot carry provenance.

The consequence appears at the moment of use. A generalised linear model in
\textsf{R} will return a fitted probability for a covariate vector ten standard
deviations outside anything it was trained on, and the number is
indistinguishable in type, class and printed appearance from one it has ample
evidence for. Extrapolation does not announce itself as uncertainty; it
announces itself as confidence, because the fitted surface is smooth and the
arithmetic is untroubled.

We propose that a prediction should be a triple---a value, a guarantee, and a
status---and that \emph{refusal} should be an admissible status.
Section~\ref{sec:certificate} states what the certificate attests.
Section~\ref{sec:conformal} recalls what split conformal prediction
guarantees and, more importantly, what it does not.
Section~\ref{sec:refusal-sec} gives
the paper's central observation: under the weighted conformal construction,
refusal is not an added safeguard but the finite-sample expression of an
infinite quantile. Section~\ref{sec:thresholds} turns to a separate and more
mundane failure,
namely that the thresholds practitioners use are compared against statistics
whose null behaviour they have not checked, and gives four measured instances.
Section~\ref{sec:impl} describes the implementation.

\section{Related work}

\paragraph{Selective prediction.} Allowing a classifier to abstain is an old
idea, introduced by \citet{chow1970} as the reject option and developed into
selective classification by \citet{elyaniv2010} and, for modern networks, by
\citet{geifman2017}. The resemblance to what we propose is close enough that
the difference must be stated precisely, because it is the difference that
carries the argument.

Selective prediction abstains on \emph{confidence}: a score derived from the
model, thresholded to trade coverage against selective risk. The rejected
region is chosen to optimise that trade-off, and the threshold is a free
parameter fixed by the practitioner's tolerance for error.

We abstain on \emph{validity}. The criterion is not that the model is unsure
but that the guarantee does not apply---Proposition~\ref{prop:refusal} makes
the condition explicit, and it involves the calibration sample and the test
point's density ratio, not the model's output at all. A confident model and an
unconfident one are refused alike if the weighted calibration mass will not
reach $1-\alpha$. There is no free parameter beyond $\alpha$ itself, which is
already the miscoverage level.

The two criteria are complementary rather than competing. A practitioner may
wish to abstain on low confidence \emph{as well}; our claim is only that
abstaining where the guarantee is void is not a tuning decision.

\paragraph{Documentation artefacts.} Model cards \citep{mitchell2019} and
datasheets \citep{gebru2021} established that a model or dataset should ship
with a structured account of its provenance and intended use. We take that
premise and add a mechanism: the account is generated from a certificate bound
to the artefact by hash, so it cannot be written by hand, cannot be edited
independently of the model, and cannot silently go stale. The documentation
literature specifies what should be recorded; the gap we address is that
nothing enforces the recording or detects its divergence.

\paragraph{Leakage.} \citet{kapoor2023} survey leakage across published
machine-learning-based science and find it pervasive and reproducible, which
is the empirical case for screening at fit time rather than on request. Our
contribution here is not a better leakage detector---the screens we implement
are deliberately shallow---but the observation in Section~\ref{sec:thresholds} that the
conventional screening statistic is unusable on rare outcomes, and the
insistence that a screen which is easy to skip will be skipped.

\paragraph{Shift detection.} \citet{rabanser2019} compare methods for
detecting dataset shift and recommend dimensionality reduction followed by
two-sample testing, which is the design we adopt. Our addition is procedural:
the test statistic must be calibrated against its own null in the setting at
hand, and Section~\ref{sec:thresholds} measures what happens when it is not.

\paragraph{Conformal prediction under shift.} Beyond the covariate-shift
weighting of \citet{tibshirani2019}, \citet{gibbs2021} give an online scheme
that adapts the miscoverage level from observed coverage, which addresses
arbitrary shift but requires outcomes to arrive. That requirement is exactly
the limitation we record at the end of Section~\ref{sec:thresholds}: where labels
are available,
adaptive methods are preferable to refusal; where they are not, the density
ratio is the only handle, and it reaches covariate shift alone.

\section{What a certificate attests}
\label{sec:certificate}

\begin{definition}[Certificate]
Let $\mathcal{D}$ be a training sample, $\mathcal{S}$ a specification
consisting of a partition rule and an ordered set of checks, and $\hat f$ the
fitted model. A \emph{certificate} is the tuple
\[
  \mathcal{C} = \bigl(\,h(\mathcal{D}),\; h(\mathcal{S}),\; h(\hat f),\;
  \{(c_k, s_k, \theta_k, I_k, v_k)\}_{k=1}^{K},\; t\,\bigr),
\]
where $h$ is a collision-resistant hash, $c_k$ identifies the $k$th check,
$s_k$ is its statistic, $\theta_k$ its threshold, $I_k$ an interval estimate
for $s_k$, $v_k \in \{\textup{pass},\textup{weak},\textup{fail},
\textup{waived},\textup{untestable},\textup{info}\}$ its verdict, and $t$ the
time of issue.
\end{definition}

The content of the definition is in what it excludes. A certificate does not
assert that $\hat f$ is accurate, that $\mathcal{D}$ was appropriate, or that
the checks were the right ones. It asserts that these checks, with these
thresholds, were run on data with this hash and produced these verdicts, and
it binds that assertion to $\hat f$ so that the two cannot be separated without
detection.

\begin{remark}
$h(\mathcal{S})$ must be computed from a structural summary of the
specification rather than from its serialisation. Checks are closures, and in
\textsf{R} a closure's serialised representation changes once the just-in-time
compiler attaches byte code to it, so a specification hashed before and after
use yields different digests. This is a small point with a large consequence:
an identifier that varies across runs of the same pipeline cannot support an
audit trail.
\end{remark}

\section{What split conformal prediction guarantees}
\label{sec:conformal}

Let $(X_i, Y_i)_{i=1}^{n+1}$ be exchangeable, let $\hat\mu$ be fitted on data
independent of the calibration sample, and let $V_i = \mathcal{A}(X_i, Y_i)$ be
nonconformity scores, for instance $|Y_i - \hat\mu(X_i)|$ in regression or
$1 - \hat p_{Y_i}(X_i)$ in classification. With
$q = V_{(\lceil (n+1)(1-\alpha)\rceil)}$ the corresponding order statistic, the
set $\hat C(x) = \{y : \mathcal{A}(x,y) \le q\}$ satisfies
\begin{equation}
  1 - \alpha \;\le\; \mathbb{P}\bigl(Y_{n+1} \in \hat C(X_{n+1})\bigr)
  \;\le\; 1 - \alpha + \tfrac{1}{n+1},
  \label{eq:split}
\end{equation}
the upper bound holding when the scores are almost surely distinct
\citep{vovk2005,lei2018}.

Four properties of \eqref{eq:split} are routinely over-read.

\paragraph{The probability is marginal.} It is taken over the joint law of the
calibration sample and the test point. It is not a statement about any
particular $x$, and coverage may be far below $1-\alpha$ on an identifiable
subpopulation while holding overall.

\paragraph{Coverage is achieved by construction.} Equation~\eqref{eq:split}
holds for \emph{any} $\hat\mu$, including a constant. It is therefore not a
diagnostic: a model that has been handed a feature which restates the outcome
attains the same coverage as an honest one. In a worked example, a leaking
model is better calibrated by a factor of thirty and produces singleton label
sets where the honest model returns the uninformative full set, while the two
are indistinguishable on coverage. Sharpness, not coverage, is the quantity
that separates them, and only the pair is informative.

\paragraph{The guarantee is finite-sample and random.} Empirical coverage on a
single test set scatters around $1-\alpha$ with standard error of order
$n_{\text{test}}^{-1/2}$. A measured shortfall is therefore not evidence of
failure, which is the motivation for interval-based verdicts in Section~\ref{sec:thresholds}.

\paragraph{Exchangeability is the whole assumption.} If the calibration and
test samples are not exchangeable, \eqref{eq:split} does not degrade
gracefully; it is void. Nothing in the construction signals this, which is the
subject of the next section.

\section{Refusal is an infinite quantile}
\label{sec:refusal-sec}

Suppose covariate shift: training and calibration data are drawn from
$P = P_X \times P_{Y|X}$ and test data from $Q = Q_X \times P_{Y|X}$, with
$Q_X \ll P_X$ and likelihood ratio $w = \mathrm{d}Q_X/\mathrm{d}P_X$.
\citet{tibshirani2019} show that exchangeability can be restored by weighting.
Define, for a test point $x$,
\begin{equation}
  p_i^w(x) = \frac{w(X_i)}{\sum_{j=1}^{n} w(X_j) + w(x)},
  \qquad
  p_\infty^w(x) = \frac{w(x)}{\sum_{j=1}^{n} w(X_j) + w(x)},
  \label{eq:weights}
\end{equation}
and let $\hat C(x) = \{ y : \mathcal{A}(x,y) \le Q_{1-\alpha}(x)\}$ where
$Q_{1-\alpha}(x)$ is the $(1-\alpha)$ quantile of the weighted measure
$\sum_i p_i^w(x)\,\delta_{V_i} + p_\infty^w(x)\,\delta_{+\infty}$. Then
$\mathbb{P}_{Q}(Y \in \hat C(X)) \ge 1-\alpha$.

The atom at $+\infty$ is not a technical convenience. It is the construction
declining to extrapolate, and it becomes decisive exactly when the test point
is unlike the calibration sample.

\begin{proposition}[Refusal condition]
\label{prop:refusal}
With the weights \eqref{eq:weights}, $Q_{1-\alpha}(x) = +\infty$---and hence
$\hat C(x)$ is the whole outcome space---if and only if
\begin{equation}
  p_\infty^w(x) > \alpha,
  \qquad\text{equivalently}\qquad
  w(x) > \frac{1-\alpha}{\alpha}\,\sum_{j=1}^{n} w(X_j).
  \label{eq:refusal}
\end{equation}
\end{proposition}

\begin{proof}
The finite atoms carry total mass $1 - p_\infty^w(x)$. The $(1-\alpha)$
quantile is finite precisely when the finite part of the measure attains mass
$1-\alpha$, that is when $1 - p_\infty^w(x) \ge 1-\alpha$. Its negation is
$p_\infty^w(x) > \alpha$; substituting \eqref{eq:weights} and rearranging gives
the stated bound on $w(x)$.
\end{proof}

Proposition~\ref{prop:refusal} is elementary, and that is the point. It says
that the guarantee itself specifies when it has nothing to offer, and the
condition is a comparison between the test point's own weight and the total
weight of the calibration sample. A software system has three options when
\eqref{eq:refusal} holds: return the trivial set, return a narrower set with no
justification, or refuse. The second is what an unweighted implementation does
silently, and it is the only one of the three that is wrong. The first and
third are the same claim, differing in whether the caller is told.

We therefore take refusal to be the correct return, with the qualification that
it must be accompanied by a reason, and that overriding it must be recorded.
The practical support condition used in implementation---refusing rows with a
feature outside the training range, or an unseen categorical level---is a
computable sufficient condition for $w(x)$ to be large, and is applied
in addition to \eqref{eq:refusal} rather than instead of it.

\begin{remark}[What weighting does not buy]
The guarantee assumes $w$ is known. In practice it is estimated, typically by
a discriminative classifier, so coverage is approximate and degrades as the
estimate does. More importantly the construction addresses covariate shift
only: if $P_{Y|X}$ has itself changed, reweighting is not merely ineffective
but actively misleading, because the returned intervals carry the appearance
of having been adjusted. Under strong shift the weighted calibration sample
has small effective size and the quantile is driven by few observations, which
in our experiments produces conservative over-coverage well before the method
fails outright.
\end{remark}

\section{Thresholds without nulls}
\label{sec:thresholds}

The second half of our argument concerns a failure that requires no theory at
all. Decisions about whether a model may be deployed are made by comparing a
statistic to a threshold. The thresholds in circulation are conventions, and
conventions are calibrated---if at all---against a null the practitioner has
not checked in their own setting. We give four instances, each measured, each
repaired by supplying the missing null.

\paragraph{A statistic whose noise floor depends on sample size.} The
population stability index is conventionally flagged above $0.2$. With ten
bins and a batch of 50 rows drawn from the training distribution \emph{itself},
its median is approximately $0.2$. The threshold therefore flags roughly 80\%
of unshifted batches at the smallest batch size the software permits. Comparing
instead against a null simulated at the observed batch size, by resampling from
the stored training frequencies, reduces this to 6\% while leaving detection of
a genuine shift unchanged (Table~\ref{tab:nulls}).

\paragraph{A statistic optimistic in sample.} A classifier two-sample test
\citep{lopezpaz2017} detects changes in dependence structure that per-feature
statistics cannot see. Its area under the curve, computed in sample, reads
approximately $0.53$ on two samples drawn from the same distribution, because a
flexible classifier can separate identical samples it is allowed to both fit
and score. Cross-fitting returns $0.485$, and because the scores are then
independent of the labels under the null, the $p$-value follows the
Mann--Whitney approximation without a permutation loop.

\paragraph{A statistic that inherits the class balance.} Screening each feature
for its ability to reproduce the outcome is conventionally done by accuracy. A
majority-class rule already attains the base rate, so on a rare outcome every
feature scores near the threshold. On a portfolio of 67,856 motor policies with
a 6.8\% claim rate, all seven rating factors scored $0.9318$ against a no-claim
rate of $0.9319$. At a 2.6\% event rate a pure-noise feature clears a $0.95$
threshold and the screen refuses every model. Scoring by area under the curve,
or mean per-class recall for multiclass outcomes, removes the dependence.

\paragraph{A point estimate treated as exact.} Comparing $\hat s$ to $\theta$
ignores that $\hat s$ is estimated. We attach a bootstrap interval $I$ and
return \textup{fail} only when $I$ lies wholly beyond $\theta$,
\textup{pass} when it lies wholly within, and \textup{weak} otherwise. The
third verdict is the substantive one: it records that the evidence did not
settle the question, rather than resolving it by the accident of which side of
$\theta$ the point estimate fell.

\begin{table}[ht]
\centering
\caption{Measured effect of supplying a null. Each row reports the rate at
which a conventional threshold fires when it should not. Simulated features
are Gaussian; see Table~\ref{tab:real} for real portfolios.}
\label{tab:nulls}
\begin{tabular}{llrr}
\toprule
Statistic & Condition & Convention & With null \\
\midrule
PSI            & unshifted, 50 rows        & 80\%   & 6\%   \\
PSI            & unshifted, 60 rows        & 69\%   & 6\%   \\
C2ST AUC       & identical samples         & 0.53   & 0.485 \\
Proxy accuracy & pure noise, 2.6\% events  & 0.974  & 0.53  \\
\bottomrule
\end{tabular}
\end{table}

\subsection{The same thresholds on real portfolios}

Simulated covariates are a weak test of a binned statistic, because real
rating factors are skewed, discrete and often categorical. We repeated both
measurements on four insurance portfolios of very different size and class
balance, drawing batches from each portfolio itself so that every flag is
false by construction (Table~\ref{tab:real}).

Both effects are larger on real data than in simulation. The conventional PSI
threshold fires on between 52\% and 92\% of unshifted 50-row batches, against
80\% for Gaussian covariates, and the simulated null brings every portfolio to
between 4\% and 8\%. At 200 rows the distinction has all but disappeared,
which is consistent with the noise floor being a small-sample effect rather
than a property of the data.

The accuracy result generalises exactly. A pure-noise feature scores the base
rate to three decimal places on all four portfolios, across base rates from
0.53 to 0.99, while the area under the curve stays at chance. On
\texttt{dataOhlsson}, a portfolio of 64,548 policies with a 1\% claim rate,
accuracy scores 0.990 on noise and so clears a 0.95 threshold: a second real
dataset on which the conventional screen would have refused every model.

\begin{table}[ht]
\centering
\caption{Four insurance portfolios. Left: score given to an added pure-noise
feature. Right: share of unshifted 50-row batches flagged, four features per
portfolio.}
\label{tab:real}
\begin{tabular}{lrrrrrr}
\toprule
& & \multicolumn{3}{c}{Noise feature} & \multicolumn{2}{c}{PSI, 50 rows} \\
\cmidrule(lr){3-5}\cmidrule(lr){6-7}
Portfolio & $n$ & base rate & accuracy & AUC & fixed 0.2 & null \\
\midrule
\texttt{dataCar}       & 67{,}856 & 0.932 & 0.932 & 0.506 & 92\% & 4\% \\
\texttt{AutoBi}        &  1{,}144 & 0.527 & 0.527 & 0.504 & 52\% & 4\% \\
\texttt{SingaporeAuto} &  7{,}483 & 0.935 & 0.935 & 0.515 & 52\% & 8\% \\
\texttt{dataOhlsson}   & 64{,}548 & 0.990 & 0.990 & 0.512 & 82\% & 4\% \\
\bottomrule
\end{tabular}
\end{table}

The four instances share a structure. In each, a threshold that is meaningful
as an effect size is applied to an estimator whose null value is not where the
threshold assumes. The repair is the same in each: establish the null
behaviour in the setting at hand, and require the statistic to clear both the
null and the effect size before acting. We suggest this pairing as the default
discipline for automated model checks, and note that it is orthogonal to which
statistics are chosen.

\begin{remark}[A limit that none of this addresses]
None of the shift statistics inspects the outcome. They report that the input
distribution has moved, which is not the same as the model having become
worse: a large shift in a feature the model ignores is harmless, and a small
shift in a decisive one may not be. Detecting degradation requires labels,
which in deployment usually have not arrived. We regard this as the principal
open problem for the framework.
\end{remark}

\section{Implementation}
\label{sec:impl}

The contract is implemented in the \textsf{R} package \texttt{attest}. Checks
execute inside the fitting function, so no code path produces a model without
them; a blocking failure refuses to issue a certificate; a model whose
certificate is not valid cannot predict; and waiving a check requires a written
reason that is stored in the certificate and reproduced in every report
thereafter. Predictions return a value, an interval or label set, a status in
$\{\textup{valid}, \textup{flagged}, \textup{refused}\}$, and a reason.

Two details are worth recording. Reports are rendered from the certificate
rather than recomputed, including the diagnostic figures, so a model card
cannot disagree with the model it describes. And certificates may be appended
to a hash-chained ledger: each entry carries the digest of its predecessor, so
that editing or removing a record is detectable. We note explicitly that a
plain digest chain is tamper-evident against accident and single edits but not
against an adversary who rewrites the chain, since the digest is recomputable;
sealing entries with a keyed hash closes that gap and moves the guarantee onto
the key.

\section{Discussion}

The argument of this paper is that the scalar return type is the root of a
class of deployment failures, and that the repair is typed rather than
procedural. Adding a validation step to a pipeline is a discipline that decays;
changing what a prediction \emph{is} does not. Proposition~\ref{prop:refusal}
suggests that this is not merely a software convention: under the weighted
conformal construction the guarantee already contains a refusal outcome, and
implementations that return a number regardless are discarding information the
mathematics supplies.

We close with the observation that motivated Section~\ref{sec:thresholds}. Every defect we found
in the implementation was found by attempting to write down a claim about
it---in documentation, in a submission note, or in this paper---and then
checking the claim. None was found by writing more code. If that generalises,
it is an argument for the position advanced here: the obligation to state what
was checked is itself a mechanism for discovering that it was not.

\section*{Acknowledgements}

Software available at \url{https://github.com/DiogoRibeiro7/attest}.

\bibliographystyle{plainnat}
\bibliography{paper,methods}

\end{document}
